\chapter{Reproducibility, remaining validation, and conclusion}
\label{ch:conclusion}\label{sec:conclusion}
\chapterpromise{A final account of the established statistical results, the executed empirical results, the unresolved implementation work, and the scope of the generalized method.}

\section{Established statistical results}
\BayesBreak is a finite Bayesian calculation over ordered partitions. For regular
exponential-family segment models with proper conjugate priors, exact segment marginal likelihoods
and posterior moments are obtained from sufficient statistics. Under a factorized partition prior,
sum-product dynamic programming yields fixed-count marginal likelihoods, $p(k\mid y)$,
changepoint marginals, changepoint-event probabilities, and Bayes curves. A separate max-sum
recursion yields the joint MAP partition. For aligned conditionally independent sequences with a
common partition, sequence-specific segment marginal likelihoods multiply. The finite latent-group
criterion has monotone alternating updates under exact maximization of its Jensen minorizer. For
nonconjugate segment models, the posterior-odds bounds are conditional on a uniform error bound for
every admissible segment score.

\section{Interpretation of the executed results}
The archive contains executed synthetic recovery, boundary-calibration, observation-model,
latent-group, nonconjugate-approximation, and finite-range timing studies, together with four
real-data fits. The real-data analyses are descriptive because independent changepoint labels are
not archived. Comparator scores that use the \BayesBreak joint-MAP changepoints as the reference
measure agreement rather than external accuracy. Unsupported computations are excluded from comparator and predictive conclusions.

\section{Implementation and validation still required}
The implementation must align known Gaussian precision, Poisson exposure and predictive
observation-family semantics with the likelihoods given in the mathematical development.
Prospective validation must cover independent exactness oracles, approximation convergence,
matched-coordinate baselines, posterior calibration and held-out predictive scores. The
array-CGH weighted fit requires regeneration after this correction; segment-count saturation
in CGH and methylation requires additional sensitivity analysis. The required experiment
designs are collected in Appendix~\ref{app:experiments}, with exact machine-readable
configurations in the repository's current experimental specification.

\section{Conclusion}
The main contribution of \BayesBreak is a generalized hierarchical Bayesian segmentation method
that combines exponential-family segment integration with exact posterior computation over
ordered partitions and extends that construction to irregular designs, multiple aligned sequences,
known groups, and finite latent groups. The generality lies in the observation-model independence of
the partition recursion once mathematically valid segment marginal likelihoods are supplied. The
hierarchical structure lies in explicit probability models for shared or group-specific
changepoints and in the separately stated latent-group optimization criterion.

This formulation is grounded in conjugate analysis, product-partition models, Bayesian
multiple-changepoint recursions, hierarchical modeling, and dynamic programming. Its limitations
are equally mathematical: exactness depends on finite segment integrals and local factorization;
shared-changepoint pooling depends on alignment and conditional independence; nonconjugate results
depend on numerical error control; and empirical conclusions are limited by the available
references, replications, and comparator definitions. Stating those conditions precisely preserves
the breadth of the generalized method without replacing statistical claims by branding language.
